\chapter[Chapter title in the TOC]{CHAPTER TITLE}


\section{Numbered Section: Sample body with common functions used}
\textit{This} thesis template --- unlike most others --- is written in LaTex\footnote{Which is far superior to MS Word.}. The Maynooth University History Department's style requirements are fully encoded in this template. You'll need to populate \texttt{Metadata.tex}, \texttt{Glossary.tex} and \texttt{Bibliography\_Master.bib}\index{Bibliography|see{VS Code}} yourself. Hyphens are - (well-known punctuation), en dashes are -- (human--machine teaming) and em dashes are --- (used for asides). 

Calling glossary items in this way allows differentiation between simple acronyms and full glossary entries: \gls{ied} Vs glossary entry.

A table might look like this:

\begin{table}[H]
	\centering
	\begin{tabular}{l r}
		\toprule
		\textbf{Component} & \textbf{Cost (€)} \\
		\midrule
		13'' drone & €742 \\
		Fibre spool & €38.3 $\text{km}^{-1}$ \\
		1 kg explosive & €50 $\text{kg}^{-1}$ \\
		2 x Electric detonator & €100  $\text{det}^{-1}$ \\
		\midrule
		Total per drone & 10 km = \textbf{€1{,}375} \& 70 km = \textbf{€3{,}673}  \\
		\bottomrule
	\end{tabular}
	\caption[Estimated fibre-optic tethered FPV-OWAD cost]{Estimated Fibre-Optic FPV-OWAD Cost (exclusive of VAT)} \label{tableref}
\end{table}

Referring to a label is done like this: Table \ref{tableref}.


\begin{figure}[htbp]
	\centering
	\includegraphics[width=0.75\textwidth]{chapters/cbrn_grief.png}
	\caption{CBRN EOD grief}
	\label{fig:grief}
\end{figure}


\begin{itemize}
	\item Posscite: \posscite{tollast_2025} does the apostrophe. 
	\item Textcite: \textcite{watling_2023} is an inline citation.
	\item Parencite: \parencite{watling_2025_panel} is a parenthetical citation.
	\item Referencing pages: Maynooth don't use pp. The colon is coded in the template, \textcite[5]{watling_2023} and \parencite[7--8]{DPR_2024}.
	\item Multiple citations is simple: \textcite{watling_2025_panel,watling_2023,Hall_2021_camouflage_platform}. Maynooth order them chronologically in the text and alphabetically in the bibliography. The template will do all of that, and there's no need to order them in multiple citations.
	\item Page number for one entry in a multiple citation. Doesn't really work, but can do for one reference: like \parencite[99--100]{watling_2023,Hall_2021_camouflage_platform}. It's better to do \parencite[99--100]{Hall_2021_camouflage_platform} and \parencite{watling_2023}.
	\item Index: \index{EOD!CBRN}
\end{itemize}

\section*{Bibliography}
Manually populate \texttt{Bibliography\_Master.bib}. Suggest using VSCode\index{VS Code!for bibliographies} to view it rather other environments (such as JabRef). Take care to populate the fields per the template bib and your bibliography will build correctly --- with the correct formatting. Titles are sentence case (not every word is capitalised) and the author's only use surnames and initials like \textbf{Thomas, A. R.}.

\begin{tikzpicture}
	\begin{axis}[
		width=13cm, height=8.5cm,
		xmin=1, xmax=25, ymin=1, ymax=5,
		xtick={1,9,16,25},
		ytick={1,2,3,4,5},
		yticklabels={1:1,2:1,3:1,4:1,5:1},
		xlabel={Defender's relative lethality $K_d/K_a$\quad(raised by UxS and persistent ISR)},
		ylabel={Attacker:defender force ratio required, $R$},
		axis lines=left,
		every axis x label/.style={at={(ticklabel cs:0.5)}, anchor=north, font=\small},
		every axis y label/.style={at={(ticklabel cs:0.5)}, anchor=south, rotate=90, font=\small},
		tick label style={font=\small},
		ymajorgrids=true,
		grid style={gray!18},
		clip=false,
		]
		
		% Ceiling: maximum local superiority dispersed infantry can mass
		\addplot[ink, dashed, line width=1pt, domain=1:25, samples=2] {3};
		
		% Required ratio R = sqrt(Kd/Ka)  (Lanchester square-law break-even)
		\addplot[amber, line width=1.6pt, domain=1:25, samples=200, name path=curve] {sqrt(x)};
		
		% Reference path along the ceiling for the shaded gap
		\path[name path=ceil] (axis cs:9,3) -- (axis cs:25,3);
		
		% Shade the gap: requirement that exceeds the infantry ceiling (x >= 9)
		\addplot[amberfill, opacity=0.18] fill between[of=curve and ceil, soft clip={domain=9:25}];
		
		% Classic 3:1 anchor at Kd/Ka = 9
		\addplot[only marks, mark=*, coral, mark size=2.6pt] coordinates {(9,3)};
		
	% Double-headed arrow spanning the gap at Kd/Ka = 14
	\draw[coral, line width=1pt, {Stealth}-{Stealth}]
	(axis cs:14,3.742) -- (axis cs:14,3);
	
	% Leader lines (dashed connectors from floating labels to referents)
	\draw[gray!60, dashed, line width=0.4pt] (axis cs:9,3) -- (axis cs:9.9,2.52);
	\draw[gray!60, dashed, line width=0.4pt] (axis cs:3.45,1.58) -- (axis cs:2.55,1.60);
	
	% Annotations
	\node[ink, anchor=west, font=\small] at (axis cs:1.4,3.2)
	{Mass dispersed infantry can generate};
	\node[coral, anchor=west, font=\small] at (axis cs:9.9,2.46)
	{Classic 3:1 rule ($K_d/K_a=9$)};
	\node[amber, anchor=west, font=\small] at (axis cs:14.5,3.42)
	{the gap --- closed only by combined arms};
	\node[amber, anchor=west, font=\small] at (axis cs:3.5,1.52)
	{Breakeven force ratio $F_{breakeven}=\sqrt{K_d/K_a}$};
		
	\end{axis}
\end{tikzpicture}
